% GENERATED FILE. Edit canonical lessons, then run npm run generate:books.
% canonical-source-hash: fnv1a64:ae1c1d95d2e4fcf6
% canonical-lessons: TA-C412-ataiyalam, TA-C412-valitavaru, TA-C412-pintanku, TA-C412-munneru, TA-C412-vilaku

\chapter{Losing the Way}
\label{ch:ta-losing-the-way}
\hlchaptermodality{412}

\begin{chapteropening}
\textbf{By the end of this chapter:} \emph{I can talk about landmarks and losing the way, and falling behind, moving forward and stepping aside.}

Complete the last lesson of chapter 412: I can talk about landmarks and losing the way, and falling behind, moving forward and stepping aside.
\end{chapteropening}

\section[\texorpdfstring{\ta{அடையாளம்} (aṭaiyāḷam)}{அடையாளம் (aṭaiyāḷam)}]{\ta{அடையாளம்} (aṭaiyāḷam) — a landmark; a sign, a mark}

\label{lesson:TA-C412-ataiyalam}



\begin{quote}
Before the new one: say the Tamil for to scrape, then the Tamil for to bump.
\end{quote}

\subsection*{You'll want to know: \ta{அடையாளம்}}
\textbf{\ta{அடையாளம்}} — \emph{aṭaiyāḷam} — \textquotedblleft{}a landmark; a sign, a mark\textquotedblright{}.

\begin{quote}
\ta{பெரிய} \ta{கோயில்தான்} \ta{எங்கள்} \ta{தெருவின்} \ta{அடையாளம்}.
\end{quote}

\begin{quote}
\emph{periya kōyiltāṉ eṅkaḷ teruviṉ aṭaiyāḷam.}
\end{quote}

\begin{quote}
The big temple is the landmark of our street.
\end{quote}

\ta{அடையாள} \ta{அட்டை} is the identity card; it uses the joining form \ta{அடையாள}.

\begin{cousinweb}
\textbf{\ta{அடையாளம்}} (\emph{aṭaiyāḷam}) is native Dravidian (DEDR 89), \textquotedblleft{}a mark, a sign, an emblem\textquotedblright{}.
\end{cousinweb}

\begin{grammarlens}[title={What you've built}]
One more word for getting around.
\end{grammarlens}

\subsection*{Your turn}
\begin{itemize}
\raggedright
  \item \emph{Say it:} \emph{aṭaiyāḷam}
  \item \emph{Say it:} \emph{aṭaiyāḷam}, once more
  \item \emph{Say it:} \emph{periya kōyiltāṉ eṅkaḷ teruviṉ aṭaiyāḷam.} — The big temple is the landmark of our street.
  \item \emph{Recall:} say the Tamil for to scrape, then the Tamil for to bump, then say \emph{aṭaiyāḷam} again
\end{itemize}

\begin{tcolorbox}[breakable,colback=teal!4,colframe=teal!35!black,title={Before you move on}]
What does \ta{அடையாளம்} mean? (\textquotedblleft{}A landmark; a sign, a mark\textquotedblright{}.) Say it once more.
\end{tcolorbox}
\section[\texorpdfstring{\ta{வழிதவறு} (vaḻitavaṟu)}{வழிதவறு (vaḻitavaṟu)}]{\ta{வழிதவறு} (vaḻitavaṟu) — to lose one's way}

\label{lesson:TA-C412-valitavaru}



\begin{quote}
Before the new one: say the Tamil for to bump, then the Tamil for a landmark.
\end{quote}

\subsection*{You'll want to know: \ta{வழிதவறு}}
\textbf{\ta{வழிதவறு}} — \emph{vaḻitavaṟu} — \textquotedblleft{}to lose one's way\textquotedblright{}.

\begin{quote}
\ta{காட்டில்} \ta{நாங்கள்} \ta{வழிதவறினோம்}.
\end{quote}

\begin{quote}
\emph{kāṭṭil nāṅkaḷ vaḻitavaṟiṉōm.}
\end{quote}

\begin{quote}
We lost our way in the forest.
\end{quote}

Past (I …): \ta{வழிதவறினேன்}, \emph{vaḻitavaṟiṉēṉ}.

\ta{வழி}, the way, + \ta{தவறு}, to slip or miss.

\begin{cousinweb}
\textbf{\ta{வழிதவறு}} (\emph{vaḻitavaṟu}) is built from \ta{வழி}, \textquotedblleft{}a way\textquotedblright{} (DEDR 5297), and the verb \ta{தவறு}, \textquotedblleft{}to slip, to miss\textquotedblright{} (as in \ta{தவறவிடு}, to miss).
\end{cousinweb}

\begin{grammarlens}[title={What you've built}]
One more word for getting around.
\end{grammarlens}

\subsection*{Your turn}
\begin{itemize}
\raggedright
  \item \emph{Say it:} \emph{vaḻitavaṟu}
  \item \emph{Say it:} \emph{vaḻitavaṟu}, once more
  \item \emph{Say it:} \emph{kāṭṭil nāṅkaḷ vaḻitavaṟiṉōm.} — We lost our way in the forest.
  \item \emph{Recall:} say the Tamil for to bump, then the Tamil for a landmark, then say \emph{vaḻitavaṟu} again
\end{itemize}

\begin{tcolorbox}[breakable,colback=teal!4,colframe=teal!35!black,title={Before you move on}]
What does \ta{வழிதவறு} mean? (\textquotedblleft{}To lose one's way\textquotedblright{}.) Say it once more.
\end{tcolorbox}
\section[\texorpdfstring{\ta{பின்தங்கு} (piṉtaṅku)}{பின்தங்கு (piṉtaṅku)}]{\ta{பின்தங்கு} (piṉtaṅku) — to lag behind; to fall behind}

\label{lesson:TA-C412-pintanku}



\begin{quote}
Before the new one: say the Tamil for a landmark, then the Tamil for to lose one's way.
\end{quote}

\subsection*{You'll want to know: \ta{பின்தங்கு}}
\textbf{\ta{பின்தங்கு}} — \emph{piṉtaṅku} — \textquotedblleft{}to lag behind; to fall behind\textquotedblright{}.

\begin{quote}
\ta{மலை} \ta{ஏறும்போது} \ta{நான்} \ta{மற்றவர்களைவிடப்} \ta{பின்தங்கினேன்}.
\end{quote}

\begin{quote}
\emph{malai ēṟumpōtu nāṉ maṟṟavarkaḷaiviṭap piṉtaṅkiṉēṉ.}
\end{quote}

\begin{quote}
While climbing the hill I fell behind the others.
\end{quote}

Past (I …): \ta{பின்தங்கினேன்}, \emph{piṉtaṅkiṉēṉ}.

\ta{பின்}, behind, + \ta{தங்கு}, to stay.

\begin{cousinweb}
\textbf{\ta{பின்தங்கு}} (\emph{piṉtaṅku}) is built from \ta{பின்}, \textquotedblleft{}the back, behind\textquotedblright{} (DEDR 4205), and \ta{தங்கு}, \textquotedblleft{}to stay\textquotedblright{} (DEDR 3014).
\end{cousinweb}

\begin{grammarlens}[title={What you've built}]
One more word for getting around.
\end{grammarlens}

\subsection*{Your turn}
\begin{itemize}
\raggedright
  \item \emph{Say it:} \emph{piṉtaṅku}
  \item \emph{Say it:} \emph{piṉtaṅku}, once more
  \item \emph{Say it:} \emph{malai ēṟumpōtu nāṉ maṟṟavarkaḷaiviṭap piṉtaṅkiṉēṉ.} — While climbing the hill I fell behind the others.
  \item \emph{Recall:} say the Tamil for a landmark, then the Tamil for to lose one's way, then say \emph{piṉtaṅku} again
\end{itemize}

\begin{tcolorbox}[breakable,colback=teal!4,colframe=teal!35!black,title={Before you move on}]
What does \ta{பின்தங்கு} mean? (\textquotedblleft{}To lag behind; to fall behind\textquotedblright{}.) Say it once more.
\end{tcolorbox}
\section[\texorpdfstring{\ta{முன்னேறு} (muṉṉēṟu)}{முன்னேறு (muṉṉēṟu)}]{\ta{முன்னேறு} (muṉṉēṟu) — to move forward; to advance}

\label{lesson:TA-C412-munneru}



\begin{quote}
Before the new one: say the Tamil for to lose one's way, then the Tamil for to lag behind.
\end{quote}

\subsection*{You'll want to know: \ta{முன்னேறு}}
\textbf{\ta{முன்னேறு}} — \emph{muṉṉēṟu} — \textquotedblleft{}to move forward; to advance\textquotedblright{}.

\begin{quote}
\ta{நெரிசலில்} \ta{வண்டிகள்} \ta{மெதுவாக} \ta{முன்னேறின}.
\end{quote}

\begin{quote}
\emph{nericalil vaṇṭikaḷ metuvāka muṉṉēṟiṉa.}
\end{quote}

\begin{quote}
In the traffic jam the vehicles moved forward slowly.
\end{quote}

Past (I …): \ta{முன்னேறினேன்}, \emph{muṉṉēṟiṉēṉ}.

Also to make progress in life: \ta{அவர்} \ta{வேலையில்} \ta{முன்னேறினார்}.

\begin{cousinweb}
\textbf{\ta{முன்னேறு}} (\emph{muṉṉēṟu}) is built from \ta{முன்}, \textquotedblleft{}the front\textquotedblright{} (DEDR 5020), and \ta{ஏறு}, \textquotedblleft{}to climb, to rise\textquotedblright{} (DEDR 916).
\end{cousinweb}

\begin{grammarlens}[title={What you've built}]
One more word for getting around.
\end{grammarlens}

\subsection*{Your turn}
\begin{itemize}
\raggedright
  \item \emph{Say it:} \emph{muṉṉēṟu}
  \item \emph{Say it:} \emph{muṉṉēṟu}, once more
  \item \emph{Say it:} \emph{nericalil vaṇṭikaḷ metuvāka muṉṉēṟiṉa.} — In the traffic jam the vehicles moved forward slowly.
  \item \emph{Recall:} say the Tamil for to lose one's way, then the Tamil for to lag behind, then say \emph{muṉṉēṟu} again
\end{itemize}

\begin{tcolorbox}[breakable,colback=teal!4,colframe=teal!35!black,title={Before you move on}]
What does \ta{முன்னேறு} mean? (\textquotedblleft{}To move forward; to advance\textquotedblright{}.) Say it once more.
\end{tcolorbox}
\section[\texorpdfstring{\ta{விலகு} (vilaku)}{விலகு (vilaku)}]{\ta{விலகு} (vilaku) — to move away; to step aside}

\label{lesson:TA-C412-vilaku}



\begin{quote}
Before the new one: say the Tamil for to lag behind, then the Tamil for to move forward.
\end{quote}

\subsection*{You'll want to know: \ta{விலகு}}
\textbf{\ta{விலகு}} — \emph{vilaku} — \textquotedblleft{}to move away; to step aside\textquotedblright{}.

\begin{quote}
\ta{வண்டி} \ta{வந்ததால்} \ta{நான்} \ta{ஓரமாக} \ta{விலகினேன்}.
\end{quote}

\begin{quote}
\emph{vaṇṭi vantatāl nāṉ ōramāka vilakiṉēṉ.}
\end{quote}

\begin{quote}
Because a vehicle came, I stepped aside.
\end{quote}

Past (I …): \ta{விலகினேன்}, \emph{vilakiṉēṉ}.

What one moves away from takes -\ta{இலிருந்து}: \ta{கூட்டத்திலிருந்து} \ta{விலகு}, to move away from a crowd.

\begin{cousinweb}
\textbf{\ta{விலகு}} (\emph{vilaku}) is native Dravidian (DEDR 5423), \textquotedblleft{}to withdraw, to step aside, to keep away\textquotedblright{}.
\end{cousinweb}

\begin{grammarlens}[title={What you've built}]
One more word for getting around.
\end{grammarlens}

\subsection*{Your turn}
\begin{itemize}
\raggedright
  \item \emph{Say it:} \emph{vilaku}
  \item \emph{Say it:} \emph{vilaku}, once more
  \item \emph{Say it:} \emph{vaṇṭi vantatāl nāṉ ōramāka vilakiṉēṉ.} — Because a vehicle came, I stepped aside.
  \item \emph{Recall:} say the Tamil for to lag behind, then the Tamil for to move forward, then say \emph{vilaku} again
\end{itemize}

\begin{tcolorbox}[breakable,colback=teal!4,colframe=teal!35!black,title={Before you move on}]
What does \ta{விலகு} mean? (\textquotedblleft{}To move away; to step aside\textquotedblright{}.) Say it once more.
\end{tcolorbox}
